\section{Introduction}
\label{s1}

\subsection{The identity and connections to integrable probability}
 Symmetric polynomials have been of much interest in various fields of mathematics such as representation theory, combinatorics etc. Among them the Schur polynomials $s_{\l}(a_1,\dots,a_N)$ occupy an important place because of their simplicity, depth and wide applications. One of their most well-known properties
 is the Cauchy identity,
\begin{align}
\label{1-1}
 \sum_{\l\in\Pe}s_{\l}(a_1,\dots,a_N)s_{\l}(b_1,\dots,b_M)
 =
 \prod_{i=1}^N\prod_{j=1}^M\frac{1}{1-a_ib_j},
\end{align}
 where $\Pe$ is the set of partitions.
 Along with a number of related summation identities, \eqref{1-1} provides important insights into various field of mathematics such as combinatorics \cite{Sagan2000symmetricgroup}, representation theory \cite{Bump2013Liegroups}, integrable probability \cite{BorodinGorin2016IntegrableProbability}, etc.

 The Cauchy identity for the Schur function \eqref{1-1} has been generalized in various different ways. In this paper, we focus on two generalizations of \eqref{1-1} in which the Schur polynomials in the left-hand side are replaced by the $q$-Whittaker functions $P_{\l}(a_1,\dots,a_N)$, $Q_{\l}(b_1,\dots,b_M)$ and the skew Schur polynomials $s_{\l/\rho}(a_1,\dots,a_N)$. These are
 \begin{align}
 \label{1-3}
 &\sum_{\mu\in\Pe}P_{\mu}(a_1,\dots,a_N)
 Q_{\mu}(b_1,\dots,b_M)
 =
 \prod_{i=1}^N\prod_{j=1}^M\frac{1}{(a_ib_j;q)_{\i}},
 \\
 \label{1-2}
 &\sum_{\substack{\l,\rho\in\Pe\\
 \rho\subset \l
 }}q^{|\rho|}s_{\l/\rho}(a_1,\dots,a_N)s_{\l/\rho}(b_1,\dots,b_M)
 =
 \frac{1}{(q;q)_\i}
 \prod_{i=1}^N\prod_{j=1}^M\frac{1}{(a_ib_j;q)_{\i}},
 \end{align}
 where $|\rho|:=\rho_1+\rho_2+\cdots$, $\rho\subset\l$ means $\rho_i\le \l_i$ for any $i=1,2,\dots$ and $(x;q)_{\i}$ is defined by~\eqref{a2-1}. The $q$-Whittaker functions $P_{\l}(a_1,\dots,a_N)$ are a special case of the Macdonald symmetric polynomials $P_{\l}(a;q;t)$ defined in \cite[Chapter~VI]{macdonald1998symmetric}, in which out of two parameters $(q,t)$ of the latter one puts $t=0$. The dual polynomial $Q_{\mu}(b_1,\dots,b_M)$ is defined as
 $
 \prod_{j=1}^M (q;q)_{\mu_j-\mu_{j+1}}^{-1}\cdot P_{\mu}(b_1,\dots,b_M)$,
 where $(x;q)_n$ is given by \eqref{a2-1d} and we set $\mu_{M+1}\equiv 0$. When $q=0$, both $P_{\l}(a_1,\dots,a_N)$ and~${Q_{\l}(a_1,\dots,a_N)}$ reduce to the Schur polynomial and \eqref{1-2} becomes \eqref{1-1}. The skew Schur polynomials \smash{$s_{\l/\rho}(a_1,\dots,a_N)$} are labeled by skew partitions $\l/\rho$. For a precise definition, see, e.g., \cite[Chapter~I]{macdonald1998symmetric}. When $q=0$, since we have $\rho=\varnothing$ due to the weight $q^{|\rho|}$, they reduce to (non-skew) Schur polynomial $s_{\l}(a_1,\dots,a_N)$ and \eqref{1-3} becomes \eqref{1-1}. Proofs of \eqref{1-3} and~\eqref{1-2} can be found in \cite[Sections~VI.4 and~I.5.28\,(a)]{macdonald1998symmetric}.

 We immediately see that the right-hand sides of \eqref{1-3} and \eqref{1-2} are exactly the same up to a factor $1/(q;q)_{\i}$. Combining this fact with \eqref{a2-8}, which rewrites $1/(q;q)_{\i}$ as a sum over partitions $\nu$, we have
 \begin{align}
 \label{1-4}
 \sum_{\mu,\nu\in\Pe}q^{|\nu|}P_{\mu}(a_1,\dots,a_N)
 Q_{\mu}(b_1,\dots,b_M)
 =
 \sum_{\substack{\l,\rho\in\Pe\\ \rho\subset\l}}q^{|\rho|}s_{\l/\rho}(a_1,\dots,a_N)s_{\l/\rho}(b_1,\dots,b_M).
 \end{align}
 This identity suggests that there may exist deeper connections between the $q$-Whittaker functions and the skew Schur functions resulting in possible refinements of \eqref{1-4}.
 Indeed, in this paper, we will prove such an identity as stated in the following theorem.

 \begin{Theorem} \label{t1}
 For any $n\in\Z_{\ge 0}$, we have
 \begin{gather}
 \sum_{\stackrel{\mu,\nu\in\Pe}{\mu_1+\nu_1\le n}}q^{|\nu|}P_{\mu}(a_1,\dots,a_N)
 Q_{\mu}(b_1,\dots,b_M)\nonumber\\
 \qquad=
 \sum_{\stackrel{\l,\rho\in\Pe}{\rho\subset\l,\,\l_1\le n}}
 q^{|\rho|}s_{\l/\rho}(a_1,\dots,a_N)s_{\l/\rho}(b_1,\dots,b_M). \label{t1-1}
 \end{gather}
 \end{Theorem}
 Summations of this type with restrictions on the length of the first row (mostly for the Schur case) have attracted attention in the theory of symmetric functions in connection with combinatorics and representation theory, see for instance \cite{Baik-Rains2001algebraic,Gessel1990Cauchysum,Rains1998IncreasingSubsequence,Rains-Warnaar2015boundedLittlewoodidentities}.
 In this paper, the proof of Theorem~\ref{t1} will be based on connections of quantities on both sides to integrable probability. In our next paper \cite{ImamuraMucciconiSasamoto2021skewRSK}, we will develop a combinatorial theory which allows a bijective proof of~\eqref{t1-1} and related identities.

 \begin{Remark} \label{rem:intro}
 In Theorem~\ref{t1}, our main identity is stated for complex variables $a_1,\dots,a_N$, $b_1,\dots,b_M$, but one easily sees that the same holds also when they are arbitrary specialization of symmetric functions. In order to establish \eqref{t1-1}, it suffices to prove it for $a_i$, $b_j$ varying in an continuous interval, which is in fact what we will prove below using integrable probabilistic arguments.
 \end{Remark}

 Let us rewrite \eqref{t1-1} in a probabilistic language. The probability measures on $\Pe$ related to~\eqref{1-3} and \eqref{1-2} are of much interest in integrable probability.
 They are defined as
 \begin{gather}
 M_{q\text{W}}(\mu)=P_{\mu}(a_1,\dots,a_N)Q_{\mu}(b_1,\dots,b_M)/Z_{q\text{W}}, \label{s13}
 \\
 M_{\text{pS}}(\l)=
 \sum_{\substack{\rho\in\Pe\\ \rho \subset\l }}q^{|\rho|}s_{\l/\rho}(a_1,\dots,a_N)s_{\l/\rho}(b_1,\dots,b_M)/Z_{\text{pS}}, \label{1-8}
 \end{gather}
 where we assumed $q \in (0,1)$ $a_i,b_j>0$, with $a_ib_j<1$ for $i=1,\dots,N$ and $j=1,\dots,M$ and the normalization constants $Z_{q\text{W}}$ and $Z_{\text{pS}}$ are given by the right-hand sides of the Cauchy identities~\eqref{1-3} and \eqref{1-2} respectively.
 $M_{q\text{W}}(\mu)$ is called the \emph{$q$-Whittaker measure}
 while $M_{\text{pS}}(\l)$ is called the \emph{periodic Schur measure}.

 The $q$-Whittaker measure was introduced in \cite{BorodinCorwin2014Mac}. It has been playing an important role in the development of integrable probability in the past decade. Its marginals are known to describe one point functions of certain particle systems discretizing the KPZ equation. On the other hand, the periodic Schur measure was introduced in \cite{Borodin2007PeriodicSchur}. It can be regarded as a model of lozenge tilings in a cylindric domain and free fermions at positive temperature \cite{BeteaBouttier2019PeriodicSchur}.

 We also introduce two random variables $\chi$ and $S$, which are independent of other random variables and respectively have distribution,
 \begin{align}
\label{t11-3}
\P(\chi=n)=\frac{q^n}{(q;q)_n}(q;q)_{\i},
\qquad
n=0,1,2,\dots
\end{align}
and for a fixed parameter $t \in \mathbb{R}_{>0}$,
 \begin{align}
 \label{1-11}
 \P(S=\ell):=\frac{t^{\ell} q^{\ell^2/2}}{(q;q)_\i
 \theta \bigl(-tq^{1/2}\bigr)},
 \qquad
\ell \in \mathbb{Z}.
\end{align}
Here the definitions of the $q$-Pochhammer symbols
$(x;q)_\i$ and $(x;q)_n$ are given by \eqref{a2-1} and~\eqref{a2-1d} respectively whereas the function $\theta(x)$ is defined below \eqref{a2-6}. Being proportional to $q^{\ell^2/2},~\ell\in\mathbb{Z}$, the distribution defined by \eqref{1-11} is sometimes called a discrete Gaussian distribution.

The random variable $\chi$ is related to the partition $\nu$ in \eqref{t1-1}.
Namely, thanks to the summation identity \eqref{a2-7}, if a partition $\nu$ is sampled with probability proportional to $q^{|\nu|}$, then its first row~$\nu_1$ follows the same law as $\chi$. On the other hand the random variable $S$ is introduced as it unveils a~determinantal structure in the periodic Schur process. This was observed by Borodin in \cite{Borodin2007PeriodicSchur}, who noticed that shifting by $S$ all entries of a~random partition $\lambda$ sampled with $M_{\mathrm{pS}}$ law produces a~determinantal point process called \emph{shift-mixed periodic Schur measure}.


With these preparations, Theorem~\ref{t1} can be restated as follows.
 \begin{Theorem} \label{thm:intro_prob}\label{t12}
 Let $\mu_1$, $\lambda_1$ be the first rows of partitions $\mu$, $\lambda$, distributed respectively according to the $q$-Whittaker measure $M_{q{\rm W}}(\mu)$ \eqref{s13} and to the periodic Schur measure $M_{{\rm pS}}(\l)$ \eqref{1-8}. Let also $\chi$ and $S$ be independent random variables distributed as \eqref{t11-3} and \eqref{1-11}, respectively.
 The following equalities hold and they are equivalent. For $n\in\Z_{\ge 0}$,
 we have
 \begin{align}
 \label{t12-1}
& {\rm (a)}\quad \P (\mu_1+\chi\le n )=
 \P (\l_1\le n ),
\\
\label{t12-3} & {\rm (b)} \quad
 \E\left[\frac{1}{\bigl(-t q^{\frac12+n-\mu_1};q\bigr)_{\i}}\right]
 =
 \P(\l_1+S\le n).
\end{align}
\end{Theorem}
 Dividing both sides of \eqref{t1-1} by the right-hand side by $Z_{\textrm{pS}}$ and noting $Z_{\textrm{pS}}=Z_{q\textrm{W}}Z_{\chi}$, where~${Z_{\chi}:=1/(q;q)_{\infty}}$ is the normalization constant appearing in \eqref{t11-3}, we get \eqref{t12-1}.
 Equivalence between \eqref{t12-1} and \eqref{t12-3} will be shown in Section~\ref{s23} although it is rather straightforward.
 The key observation is the formula
 \begin{align}
 \label{chiS}
 \P(\chi+S\le n)
 =
 \frac{1}{\bigl(-t q^{\frac12+n};q\bigr)_{\i}}.
 \end{align}
 This was already shown in \cite[Section 2]{BeteaBouttier2019PeriodicSchur}, but we will give a more direct proof in Lemma~\ref{l27} below. From \eqref{chiS}, we easily see that \eqref{t12-3} can be rewritten as{\samepage
 \[\P(\mu_1+\chi+S\le n)
 =
 \P(\l_1+S\le n),\]
 which is clearly equivalent to \eqref{t12-1}.}

 Now our goal is to prove \eqref{t12-3}. At this point we remark that, thanks to results from integrable probability \cite{Borodin2007PeriodicSchur,ImamuraSasamoto2019qTASEPtheta}, the expectation and the probability on both sides are known to be written as Fredholm determinants, see Propositions~\ref{p22} for the expectation on the left-hand side and Propositions~\ref{p25} for the probability on the right-hand side.
 Crucially, the kernels of the two Fredholm determinants turn out to be very similar: they both possess complex contour integral expressions with the integrands
 being the same and
 they differ from each other only by a choice of the contours.
 What we actually prove in this paper is that this difference of contours does not produce a difference to the value of the Fredholm determinants, see Theorem~\ref{t31}.\looseness=1


